\documentclass[10pt, a4paper]{article}

\usepackage[utf8]{inputenc}
\usepackage[russian]{babel}
\usepackage{enumitem}
\usepackage{hyperref}
\usepackage{fontawesome5}
\usepackage[
    left=1.5cm,
    right=1.5cm,
    top=1.5cm,
    bottom=1.5cm
]{geometry}
\usepackage{titlesec}

\pagenumbering{gobble}

\setlength{\parindent}{0pt}
\setlength{\parskip}{4pt}

% Компактные списки
\setlist[itemize]{
    noitemsep,
    topsep=0pt
}

% Заголовки разделов
\titleformat{\section}
    {\large\bfseries}
    {}
    {0pt}
    {}

\titlespacing*{\section}
    {0pt}
    {10pt}
    {4pt}

% Подзаголовки
\titleformat{\subsection}
    {\bfseries}
    {}
    {0pt}
    {}

\titlespacing*{\subsection}
    {0pt}
    {6pt}
    {2pt}

\hypersetup{
    colorlinks=true,
    urlcolor=black,
    linkcolor=black
}

\begin{document}

%================================================
% HEADER
%================================================

\begin{center}
    {\LARGE \textbf{Артём Шкарбан}} \\[0.5em]

    \faEnvelope\
    \href{mailto:shkarban.ab@phystech.edu}
    {shkarban.ab@phystech.edu}
    \quad
    \faTelegram\ @artemarte1978
\end{center}

%================================================
% SKILLS
%================================================

\section*{НАВЫКИ}

SystemVerilog
\quad
Verilog
\quad
FPGA
\quad
Xilinx Vivado
\quad
UVM
\quad
Xcelium
\quad
Genus
\quad
Git
\quad
Linux
\quad
Tcl
\quad
Python
\quad
MATLAB

\rule{\linewidth}{2pt}

%================================================
% EXPERIENCE
%================================================

\section*{ОПЫТ РАБОТЫ}

%------------------------------------------------
% YADRO
%------------------------------------------------

\subsection*{Инженер аппаратного обеспечения}

\textbf{YADRO, Октябрь 2025 --- настоящее время}

\textbf{Задачи}

\begin{itemize}
    \item Изучение алгоритмов POLAR, LDPC и TBCC-кодирования в рамках стандартов 5G NR.
    \item Разработка функциональных RTL-модулей на SystemVerilog для блоков канального кодирования 5G NR.
    \item Поддержка Assembly Unit (AU) на этапах проверки дизайна с использованием LINT, LEC, SYN и Litmus.
    \item Верификация разработанных блоков совместно с командой верификации с использованием UVM и regression-тестирования.
\end{itemize}

\textbf{Результаты}

\begin{itemize}
    \item На основе изученных спецификаций разработал функциональные RTL-блоки, вошедшие в состав ускорителей POLAR ENC и LDPC ENC.
    \item Дорабатывал и интегрировал функциональные модули в составе блоков канального кодирования.
    \item Исправлял функциональные и интеграционные ошибки, выявленные в ходе UVM-регрессий.
    \item Поддерживал разработанные блоки до прохождения LINT, LEC и синтеза.
\end{itemize}

\rule{\linewidth}{1pt}

%------------------------------------------------
% MIPT LAB
%------------------------------------------------

\subsection*{Инженер-программист}

\textbf{Лаборатория прикладных нанотехнологий МФТИ, Декабрь 2024 --- Октябрь 2025}

\textbf{Задачи}

\begin{itemize}
    \item Разработка IP-блоков на SystemVerilog для FPGA Intel (Altera) и Xilinx.
    \item Верификация RTL-кода.
    \item Моделирование помехоустойчивых систем в MATLAB.
\end{itemize}

\textbf{Результаты}

\begin{itemize}
    \item Доработал существующую модель BCH-декодера в MATLAB.
    \item На основе MATLAB-модели разработал RTL-модули BCH-энкодера и BCH-декодера.
    \item Провёл функциональную верификацию разработанных модулей с использованием тестбенчей.
\end{itemize}

\rule{\linewidth}{1pt}

%------------------------------------------------
% LEPTON
%------------------------------------------------

\subsection*{Инженер-программист}

\textbf{НПО «Лептон», Ноябрь 2022 --- Сентябрь 2024}

\textbf{Задачи}

\begin{itemize}
    \item Разработка IP-блоков и их реализация на FPGA Intel.
    \item Разработка Python-скриптов для взаимодействия с внешними устройствами.
\end{itemize}

\textbf{Результаты}

\begin{itemize}
    \item Реализовал интерфейсы взаимодействия FPGA с датчиками температуры и давления.
    \item Реализовал передачу данных от внешних устройств через FPGA на компьютер.
\end{itemize}

\rule{\linewidth}{2pt}

%================================================
% EDUCATION
%================================================

\section*{ОБРАЗОВАНИЕ}

\textbf{Московский физико-технический институт (МФТИ)}

\textbf{Факультет радиотехники и кибернетики (ФРТК), 2025 --- настоящее время, магистратура}

\vspace{4pt}

\textbf{Московский физико-технический институт (МФТИ)}

\textbf{Факультет аэрокосмических исследований (ФАКИ), 2020 --- 2025, бакалавриат}

\end{document}